\section{Related Work}

\paragraph{LLM-as-a-Judge and reward evaluation.}
Recent work studies whether LLMs can replace or assist human evaluators in
open-ended tasks, often by measuring agreement with human preferences or
benchmark labels. G-Eval uses LLMs with chain-of-thought and form-filling
prompts for reference-free NLG evaluation~\cite{liu2023geval}. MT-Bench and
Chatbot Arena popularized strong LLMs as pairwise
judges~\cite{zheng2023judging}, while RewardBench and related
benchmarks evaluate reward models and preference-ranking
behavior~\cite{lambert2024rewardbench}. These works primarily assess the final
decision layer: whether the judge prefers the right answer. Our work studies
an earlier layer. We ask whether the judge has generated or considered the
right evaluation dimensions before making a decision. This makes blind-spot
coverage complementary to judge accuracy, preference alignment, and
reward-model ranking metrics: those metrics can be high even when the
evaluative basis is incomplete for harder neighboring cases.
The distinction is between evaluating a judge's output and auditing the
dimension set that supports the judgment. A preference label can certify that a
decision was correct on one instance, but it does not certify that the judge
would attend to safety, constraint following, evidence grounding, or other
human-prioritized dimensions when those dimensions become decisive.
This is the blind-spot setting: the observable judgment can be acceptable while
the latent evaluation basis is missing dimensions that would matter under a
nearby response pair or a stricter human review.
A high preference-agreement score is therefore evidence about observed
decisions, not evidence that the judge represented every human-gold evaluation
dimension needed for robust evaluation.

\paragraph{Evaluation criteria for open-ended evaluation.}
Explicit criteria make subjective evaluation more structured by decomposing
quality into decision dimensions. Existing evaluation-criteria elicitation
systems usually optimize usefulness, specificity, or human preference over
criteria text. Criteria-conditioned evaluators such as Prometheus show that explicit score criteria can make
long-form evaluation more fine-grained and controllable~\cite{kim2023prometheus}.
However, these systems generally assume that the supplied criteria already
contain the right dimensions. We instead define a dimension-level coverage
target against human-gold dimensions. The goal is not to make criteria longer or
more fluent, but to measure which human evaluation dimensions are omitted and
test when dimension-level coverage-change or recovery wording is supported
under the evidence gates.
Thus the unit of analysis is not the rubric text as an artifact, but the
semantic set of evaluation dimensions it covers. This lets us separate a
supported dimension-level coverage-change statement from verbosity, paraphrase
diversity, or stylistic variation in the generated criteria.
It also changes the optimization target: the objective is not to imitate a
preferred wording style, but to cover human-prioritized dimensions under a
fixed semantic matching and validity protocol.

\paragraph{RLVR and GRPO.}
RLVR has been most successful in domains with externally checkable answers,
such as mathematics, coding, and tool-use tasks. GRPO-style optimization
further reduces dependence on learned value models by comparing grouped
rollouts under verifiable rewards~\cite{shao2024deepseekmath,guo2025deepseekr1}.
Open-ended criteria elicitation lacks exact answers, so it is not an obvious RLVR target.
BSC provides the missing bridge: the generated text remains open-ended, but the
reward is computed from semantic coverage, validity, and redundancy under a
fixed protocol. This is different from prompt tuning an evaluation-criteria
policy for surface-form optimization: the reward is tied to whether human evaluation
dimensions are covered, non-redundant, and decidable.
The resulting research question is whether verifiable rewards can supervise
semantic dimension elicitation when exact textual answers are unavailable, not
whether reinforcement learning can improve a downstream prompt in an
uncontrolled way.

\paragraph{Data contamination and proxy supervision.}
Because evaluation criteria can leak directly into training targets, data
isolation is central to this problem. BlindSpot-RL separates human hard-gold
evaluation from proxy-gold training. Human-gold RubricBench \texttt{test\_main}
is reserved for main BSC evaluation; proxy-gold criteria from RewardBench,
IFBench, WritingBench, HealthBench, and BEIR/NQ are elicited from teachers and
filtered before training. This separation is part of the scientific claim, not
just an implementation detail.
